\chapter{Obesity}
\index{Obesity}
\label{sec:Obesity}

\section{Overview}


\begin{itemize}[noitemsep]
  \item Obesity is a chronic, relapsing, multifactorial disease characterized by excessive accumulation of adipose tissue with adverse health consequences. The global prevalence has nearly tripled since 1975
  \item Over 650 million adults are obese.
\end{itemize}
\section{Causes}
Clinical consequences affect virtually every organ system: type 2 diabetes mellitus, high blood pressure, dyslipidemia, coronary artery disease, stroke, obstructive sleep apnea, non-alcoholic fatty liver disease, cholelithiasis, osteoarthritis, infertility, depression, and increased risk of multiple cancers. Diagnosis includes BMI measurement, waist circumference, and screening for comorbidities. Metabolic disorders arise from dysregulation of normal biochemical pathways. This may result from enzyme deficiencies, hormonal imbalances, or lifestyle factors that overwhelm normal homeostatic mechanisms. The underlying mechanisms involve complex interactions between genetic predisposition and environmental factors. Disruption of normal physiological processes leads to the characteristic features of the condition. Multiple contributing factors may act synergistically to trigger or worsen the condition.
\section{Risk Factors}
\begin{itemize}[noitemsep]
  \item Genetic predisposition and family history
  \item Advancing age
  \item Lifestyle factors (diet, exercise, smoking, alcohol)
  \item Environmental exposures
  \item Underlying medical conditions
  \item Medication use
\end{itemize}
\section{Signs and Symptoms}

\subsection{Common symptoms}
\begin{itemize}[noitemsep]
  \item Pain or discomfort in the affected area
  \end{itemize}

\subsection{Emergency warning signs}
\begin{itemize}[noitemsep]
  \item Severe or worsening symptoms of obesity require immediate medical evaluation
\end{itemize}
\section{Progression and Stages}
The condition may progress through different stages of severity over time. Early stages often involve mild or subtle symptoms that may be overlooked. Moderate stages involve more noticeable symptoms affecting daily function. Advanced stages may involve significant impairment and complications.

\begin{itemize}[noitemsep]
  \item It is defined by a body mass index (BMI) of 30 kg/m\textsuperscript{2} or greater in adults, classified into class I (30--34.9), class II (35--39.9), and class III (40 or higher). This condition happens because of a chronic energy imbalance---caloric intake exceeding energy expenditure---driven by complex interactions of genetic, environmental, behavioral, psychosocial, endocrine, and iatrogenic factors
  \item Adipose tissue is an active endocrine organ secreting adipokines (leptin, adiponectin, resistin, TNF$\alpha$, IL-6) contributing to a chronic low-grade inflammatory state. Treatment typically involves stepwise: lifestyle intervention (dietary modification, physical activity, behavioral therapy), pharmacotherapy (orlistat, GLP-1 receptor agonists such as semaglutide and liraglutide, naltrexone-bupropion, phentermine-topiramate), and bariatric surgery for class III obesity or class II with comorbidities.
\end{itemize}
\section{Complications}
\begin{itemize}[noitemsep]
  \item Disease progression leading to worsening symptoms
  \item Secondary infections or injuries
  \item Side effects of treatment
  \item Reduced quality of life
  \item Emotional and psychological impact
\end{itemize}
\section{Diagnosis}
\begin{itemize}[noitemsep]
  \item Comprehensive medical history and physical examination
  \item Laboratory tests to assess organ function
  \item Imaging studies to visualize affected structures
  \item Specialized testing depending on the affected system
  \item Consultation with relevant specialists
\end{itemize}
\section{Prevention}
\begin{itemize}[noitemsep]
  \item Maintain a healthy lifestyle with balanced diet and exercise
  \item Avoid known risk factors
  \item Attend regular medical check-ups for early detection
  \item Follow recommended screening guidelines
\end{itemize}
\section{Treatment Options}
Dietary modifications and nutritional counseling Pharmacotherapy to correct metabolic abnormalities Hormone replacement therapy when indicated Regular monitoring of metabolic parameters Weight management programs Exercise prescription Management of associated conditions Patient education for self-management

\begin{itemize}[noitemsep]
  \item Medications to manage symptoms and address underlying mechanisms
  \item Lifestyle modifications and self-care measures
  \item Physical therapy and rehabilitation
  \item Surgical intervention when indicated
  \item Regular monitoring and follow-up care
\end{itemize}
\section{Living with It}
\begin{itemize}[noitemsep]
  \item Follow your treatment plan as prescribed
  \item Maintain regular follow-up with your healthcare provider
  \item Make necessary lifestyle adjustments
  \item Seek support from family, friends, or support groups
  \item Stay informed about your condition
\end{itemize}
\section{Outlook}
\begin{itemize}[noitemsep]
  \item The outlook varies from person to person
  \item Sustained weight loss improves or resolves many comorbidities.
\end{itemize}
